\documentclass[12pt,a4paper]{article}
\usepackage{geometry}
\usepackage{fontspec}
\usepackage{polyglossia}
\usepackage{amsmath,amssymb}
\usepackage{graphicx}
\usepackage{float}
\usepackage{caption}
\geometry{left=30mm,right=20mm,top=20mm,bottom=20mm}
\setdefaultlanguage{russian}
\setmainfont{Times New Roman}
\setsansfont{Arial}
\raggedbottom
\setlength{\parindent}{1.25cm}
\setlength{\parskip}{0.18em}
\setlength{\abovedisplayskip}{7pt plus 2pt minus 3pt}
\setlength{\belowdisplayskip}{7pt plus 2pt minus 3pt}
\setlength{\abovedisplayshortskip}{5pt plus 2pt minus 2pt}
\setlength{\belowdisplayshortskip}{5pt plus 2pt minus 2pt}
\setlength{\textfloatsep}{9pt plus 2pt minus 2pt}
\setlength{\floatsep}{8pt plus 2pt minus 2pt}
\setlength{\intextsep}{8pt plus 2pt minus 2pt}
\captionsetup{font=small,justification=centering,singlelinecheck=false,skip=3pt}
\begin{document}

\begin{titlepage}
\thispagestyle{empty}
\noindent
\begin{minipage}[t]{0.19\textwidth}
\vspace{0pt}
\centering

\includegraphics[width=24mm]{\detokenize{assets/title_emblem.jpeg}}
\end{minipage}
\hfill
\begin{minipage}[t]{0.76\textwidth}
\vspace{0pt}
\begin{center}
{\small Министерство науки и высшего образования Российской Федерации\\Федеральное государственное автономное образовательное учреждение высшего образования\\«Московский государственный технический университет имени Н.Э. Баумана\\(национальный исследовательский университет)»\\(МГТУ им. Н.Э. Баумана)}
\end{center}
\end{minipage}

\vspace{0.35cm}
\noindent\rule{\textwidth}{1.4pt}

\vspace{0.7cm}
\begin{center}
{\bfseries ФАКУЛЬТЕТ «РОБОТОТЕХНИКА И КОМПЛЕКСНАЯ АВТОМАТИЗАЦИЯ»}\\[0.35cm]
{\bfseries КАФЕДРА РК-9 «КОМПЬЮТЕРНЫЕ СИСТЕМЫ АВТОМАТИЗАЦИИ ПРОИЗВОДСТВА»}
\end{center}

\vspace{1.65cm}
\begin{center}
{\large\bfseries Домашнее задание №1}\\[0.2cm]
{\normalsize по курсу}\\[0.2cm]
{\large\bfseries «Имитационное моделирование технологических производственных процессов»}
\end{center}

\vfill
\begin{flushright}
\begin{tabular}{@{}p{3.4cm}p{8.6cm}@{}}
Студент: & Бадретдинов Эмиль Рустэмович\\[0.15cm]
Группа: & РК9-83Б\\[0.15cm]
Преподаватель: & Берчун Юрий Валерьевич
\end{tabular}
\end{flushright}

\vfill
\begin{center}
2026 г.
\end{center}
\end{titlepage}

\section*{Задача №1. Проектирование Call-центра.}
Известно, что среднее время между звонками клиентов составляет Tc секунд, а среднее время обслуживания Ts секунд. Все потоки случайных событий считать пуассоновскими.

\noindent\textbf{1.} Рассмотреть систему без очереди. Построить графики от числа операторов: вероятности отказа вплоть до обеспечения отказов менее 1\%; математического ожидания числа занятых операторов; коэффициента загрузки операторов.

\noindent\textbf{Исходные данные:}\\
Tc = 11\\
Ts = 68

\begin{figure}[H]
\centering
\includegraphics[width=0.9\textwidth]{\detokenize{assets/task1_1__scheme.png}}
\caption*{Рисунок 1.1.1. Расчетная схема.}
\end{figure}
Здесь состояния S0, S1, ..., Sn соответствуют числу занятых операторов; λ = 0.090909 1/с, μ = 0.014706 1/с.

\noindent\textbf{Распределение состояний.} Система без очереди полностью задаётся числом занятых операторов, поэтому после распределения состояний сразу определяются загрузка операторов и вероятность отказа.

\[
a = \lambda / \mu
\]

\[
p_0 = \left(\sum_{k=0}^{n}\frac{a^k}{k!}\right)^{-1}
\]

\[
p_k = \frac{a^k}{k!}p_0,\quad k=0,\ldots,n
\]
\noindent\textbf{Число занятых операторов.} После нахождения p\_k среднее число занятых операторов и их средняя загрузка вычисляются по тому же распределению состояний.

\[
M_{\mathrm{зан}} = \sum_{k=0}^{n}k\,p_k
\]

\[
K_{\mathrm{загр}} = M_{\mathrm{зан}}/n
\]
\begin{figure}[H]
\centering
\includegraphics[width=0.8\textwidth]{\detokenize{../figures/task1_1__busy_operators_vs_operators.png}}
\caption*{Рисунок 1.1.2. Среднее число занятых операторов при изменении числа операторов.}
\end{figure}
\noindent\textbf{Вероятность отказа.} Отказ возникает только при полной занятости всех операторов, поэтому отдельного анализа требует именно крайнее состояние цепи.

\[
P_{\mathrm{отк}} = p_n
\]
\begin{figure}[H]
\centering
\includegraphics[width=0.8\textwidth]{\detokenize{../figures/task1_1__refusal_and_utilization_vs_operators.png}}
\caption*{Рисунок 1.1.3. Вероятность отказа и загрузка операторов при изменении числа операторов.}
\end{figure}
Перебор по числу операторов выполнен до первого выполнения условия P\_отк < 0.01.

Для текущего варианта минимальное число операторов равно 13, при этом вероятность отказа составляет 0.00642.

\clearpage
\noindent\textbf{2.} Рассмотреть систему с ограниченной очередью. Варьируя число операторов вплоть до 15, построить семейства графиков от числа мест в очереди: вероятности отказа; математического ожидания числа занятых операторов; коэффициента загрузки операторов; вероятности существования очереди; математического ожидания длины очереди; коэффициента занятости мест в очереди. Варьируя число мест в очереди вплоть до 15, построить семейства графиков от числа операторов для тех же метрик.

\noindent\textbf{Исходные данные:}\\
Tc = 11\\
Ts = 68

\begin{figure}[H]
\centering
\includegraphics[width=0.9\textwidth]{\detokenize{assets/task1_2__scheme.png}}
\caption*{Рисунок 1.2.1. Расчетная схема.}
\end{figure}
После заполнения всех n операторов правые состояния схемы соответствуют длине очереди от 1 до m; λ = 0.090909 1/с, μ = 0.014706 1/с.

\noindent\textbf{Распределение состояний.} После состояний с k = 0, ..., n занятыми операторами добавляются состояния очереди длины r = 1, ..., m, поэтому то же распределение используется и для потерь, и для характеристик ожидания.

\[
a = \lambda / \mu,\quad \rho_n = \lambda /(n\mu)
\]

\[
p_0 = \left(\sum_{k=0}^{n}\frac{a^k}{k!} + \frac{a^n}{n!}\sum_{r=1}^{m}\rho_n^{\,r}\right)^{-1}
\]

\[
p_k = \frac{a^k}{k!}p_0,\quad k=0,\ldots,n
\]

\[
p_{n+r} = \frac{a^n}{n!}\rho_n^{\,r}p_0,\quad r=1,\ldots,m
\]
Построены оба требуемых sweep-режима: по числу мест в очереди при фиксированном числе операторов и по числу операторов при фиксированном числе мест в очереди.

Во всех сериях использованы диапазоны по операторам 1..15 и по местам в очереди 1..15 без ручного отбора точек.

\noindent\textbf{Семейства по числу мест в очереди.} Сначала фиксируется число операторов n, а по оси абсцисс изменяется число мест в очереди m.

\noindent\textbf{Вероятность отказа.} При фиксированном n вероятность отказа определяется последним состоянием цепи с заполненной очередью.

\[
P_{\mathrm{отк}} = p_{n+m}
\]
\begin{figure}[H]
\centering
\includegraphics[width=0.8\textwidth]{\detokenize{../figures/task1_2__refusal_vs_queue__family_by_operators.png}}
\caption*{Рисунок 1.2.2. Вероятность отказа при изменении m и разных n.}
\end{figure}
\noindent\textbf{Занятые операторы и их загрузка.} По тем же точкам отдельно прослеживаются среднее число занятых операторов и относительная загрузка имеющегося ресурса обслуживания.

\[
M_{\mathrm{зан}} = \sum_{k=0}^{n+m}\min(k,n)\,p_k
\]
\begin{figure}[H]
\centering
\includegraphics[width=0.8\textwidth]{\detokenize{../figures/task1_2__busy_operators_vs_queue__family_by_operators.png}}
\caption*{Рисунок 1.2.3. Среднее число занятых операторов при изменении m и разных n.}
\end{figure}
\begin{figure}[H]
\centering
\includegraphics[width=0.8\textwidth]{\detokenize{../figures/task1_2__operators_utilization_vs_queue__family_by_operators.png}}
\caption*{Рисунок 1.2.4. Коэффициент загрузки операторов при изменении m и разных n.}
\end{figure}
\noindent\textbf{Очередь.} Следующие графики показывают, как уменьшение потерь перераспределяется в вероятность появления очереди, её среднюю длину и занятость мест.

\[
P_{\mathrm{оч}} = \sum_{k=n+1}^{n+m}p_k
\]

\[
L_{\mathrm{оч}} = \sum_{k=n+1}^{n+m}(k-n)\,p_k,\quad K_{\mathrm{мест}} = L_{\mathrm{оч}}/m
\]
\begin{figure}[H]
\centering
\includegraphics[width=0.8\textwidth]{\detokenize{../figures/task1_2__queue_exists_vs_queue__family_by_operators.png}}
\caption*{Рисунок 1.2.5. Вероятность наличия очереди при изменении m и разных n.}
\end{figure}
\begin{figure}[H]
\centering
\includegraphics[width=0.8\textwidth]{\detokenize{../figures/task1_2__queue_length_vs_queue__family_by_operators.png}}
\caption*{Рисунок 1.2.6. Средняя длина очереди при изменении m и разных n.}
\end{figure}
\begin{figure}[H]
\centering
\includegraphics[width=0.8\textwidth]{\detokenize{../figures/task1_2__queue_occupancy_vs_queue__family_by_operators.png}}
\caption*{Рисунок 1.2.7. Занятость мест очереди при изменении m и разных n.}
\end{figure}
При фиксированном n = 5 рост числа мест в очереди с m = 1 до m = 15 снижает вероятность отказа с 0.315561 до 0.19512, но одновременно переводит потерянные вызовы в ожидание: вероятность существования очереди растёт с 0.315561 до 0.978298, а её средняя длина увеличивается с 0.315561 до 11.170502.

Математическое ожидание числа занятых операторов на том же примере растёт лишь с 4.231078 до 4.975619, то есть дополнительные места прежде всего уменьшают потери, а не создают новый ресурс обслуживания.

\noindent\textbf{Семейства по числу операторов.} Затем фиксируется число мест в очереди m, а по оси абсцисс изменяется число операторов n.

\noindent\textbf{Вероятность отказа.} При фиксированном m рост числа операторов меняет как общий уровень отказов, так и скорость их убывания; используется то же выражение для P\_отк, что и в предыдущем семействе.

\begin{figure}[H]
\centering
\includegraphics[width=0.8\textwidth]{\detokenize{../figures/task1_2__refusal_vs_operators__family_by_queue.png}}
\caption*{Рисунок 1.2.8. Вероятность отказа при изменении n и разных m.}
\end{figure}
\noindent\textbf{Занятые операторы и их загрузка.} На том же sweep сравниваются те же характеристики M\_зан и загрузки операторов, но уже при изменении n.

\begin{figure}[H]
\centering
\includegraphics[width=0.8\textwidth]{\detokenize{../figures/task1_2__busy_operators_vs_operators__family_by_queue.png}}
\caption*{Рисунок 1.2.9. Среднее число занятых операторов при изменении n и разных m.}
\end{figure}
\begin{figure}[H]
\centering
\includegraphics[width=0.8\textwidth]{\detokenize{../figures/task1_2__operators_utilization_vs_operators__family_by_queue.png}}
\caption*{Рисунок 1.2.10. Коэффициент загрузки операторов при изменении n и разных m.}
\end{figure}
\noindent\textbf{Очередь.} Последний блок показывает, как увеличение n одновременно уменьшает вероятность очереди, её длину и занятость выделенных мест при тех же формулах для P\_оч, L\_оч и K\_мест.

\begin{figure}[H]
\centering
\includegraphics[width=0.8\textwidth]{\detokenize{../figures/task1_2__queue_exists_vs_operators__family_by_queue.png}}
\caption*{Рисунок 1.2.11. Вероятность наличия очереди при изменении n и разных m.}
\end{figure}
\begin{figure}[H]
\centering
\includegraphics[width=0.8\textwidth]{\detokenize{../figures/task1_2__queue_length_vs_operators__family_by_queue.png}}
\caption*{Рисунок 1.2.12. Средняя длина очереди при изменении n и разных m.}
\end{figure}
\begin{figure}[H]
\centering
\includegraphics[width=0.8\textwidth]{\detokenize{../figures/task1_2__queue_occupancy_vs_operators__family_by_queue.png}}
\caption*{Рисунок 1.2.13. Занятость мест очереди при изменении n и разных m.}
\end{figure}
При фиксированном m = 5 увеличение числа операторов с n = 5 до n = 11 снижает вероятность отказа с 0.229971 до 0.001452 и среднюю длину очереди с 2.687007 до 0.062303.

Одновременно коэффициент загрузки операторов падает с 0.952036 до 0.561168, поэтому правое семейство графиков показывает цену снижения потерь и ожидания через недогрузку части операторов.

\clearpage
\noindent\textbf{3.} Рассмотреть систему без ограничений на длину очереди. Построить графики от числа операторов вплоть до 15: математического ожидания числа занятых операторов; коэффициента загрузки операторов; вероятности существования очереди; математического ожидания длины очереди.

\noindent\textbf{Исходные данные:}\\
Tc = 11\\
Ts = 68

\begin{figure}[H]
\centering
\includegraphics[width=0.9\textwidth]{\detokenize{assets/task1_3__scheme.png}}
\caption*{Рисунок 1.3.1. Расчетная схема.}
\end{figure}
Правые состояния схемы описывают бесконечный хвост очереди после полной занятости операторов; λ = 0.090909 1/с, μ = 0.014706 1/с.

\noindent\textbf{Стационарное распределение.} Для бесконечной очереди сначала фиксируется условие стационарности ρ\_n < 1; только в устойчивых точках допускается вычисление стационарных средних характеристик.

\[
a = \lambda / \mu,\quad \rho_n = \lambda /(n\mu)
\]

\[
p_0 = \left(\sum_{k=0}^{n-1}\frac{a^k}{k!} + \frac{a^n}{n!(1-\rho_n)}\right)^{-1},\quad \rho_n < 1
\]

\[
P_{\mathrm{wait}} = \frac{a^n}{n!(1-\rho_n)}p_0
\]
Нестационарными остаются режимы n = 1, 2, 3, 4, 5, 6; для них стационарные метрики в отчёте не подставляются.

Первое стационарное значение достигается при n = 7.

\noindent\textbf{Занятые операторы.} В устойчивом режиме offered load a полностью определяет среднее число занятых операторов, а коэффициент загрузки нормирует это среднее по n.

\[
M_{\mathrm{зан}} = a,\quad K_{\mathrm{загр}} = a/n
\]
\begin{figure}[H]
\centering
\includegraphics[width=0.8\textwidth]{\detokenize{../figures/task1_3__busy_operators_vs_operators.png}}
\caption*{Рисунок 1.3.2. Среднее число занятых операторов при изменении числа операторов.}
\end{figure}
\begin{figure}[H]
\centering
\includegraphics[width=0.8\textwidth]{\detokenize{../figures/task1_3__operators_utilization_vs_operators.png}}
\caption*{Рисунок 1.3.3. Коэффициент загрузки операторов при изменении числа операторов.}
\end{figure}
\noindent\textbf{Ожидание в очереди.} После определения P\_wait стационарные характеристики очереди выражаются через ρ\_n и показывают цену приближения к границе устойчивости.

\[
P_{\mathrm{оч}} = P_{\mathrm{wait}}\rho_n
\]

\[
L_{\mathrm{оч}} = \frac{P_{\mathrm{wait}}\rho_n}{1-\rho_n}
\]
\begin{figure}[H]
\centering
\includegraphics[width=0.8\textwidth]{\detokenize{../figures/task1_3__queue_exists_vs_operators.png}}
\caption*{Рисунок 1.3.4. Вероятность наличия очереди при изменении числа операторов.}
\end{figure}
\begin{figure}[H]
\centering
\includegraphics[width=0.8\textwidth]{\detokenize{../figures/task1_3__queue_length_vs_operators.png}}
\caption*{Рисунок 1.3.5. Средняя длина очереди при изменении числа операторов.}
\end{figure}
\clearpage
\noindent\textbf{4.} Рассмотреть систему без ограничений на длину очереди, учитывающую фактор ухода клиентов из очереди при среднем приемлемом времени ожидания Tw. Построить графики от числа операторов вплоть до 15: математического ожидания числа занятых операторов; коэффициента загрузки операторов; вероятности существования очереди; математического ожидания длины очереди.

\noindent\textbf{Исходные данные:}\\
Tc = 11\\
Ts = 68\\
Tw = 105

\begin{figure}[H]
\centering
\includegraphics[width=0.9\textwidth]{\detokenize{assets/task1_4__scheme.png}}
\caption*{Рисунок 1.4.1. Расчетная схема.}
\end{figure}
После состояния полной занятости операторов уход из очереди учитывается в суммарной интенсивности обратных переходов; λ = 0.090909 1/с, μ = 0.014706 1/с, ν = 0.009524 1/с.

\noindent\textbf{Рекуррентное распределение.} В модели с уходом клиентов стационарное распределение строится рекуррентно по birth-death переходам, где после заполнения операторов обратные переходы ускоряются за счёт ухода ожидающих клиентов.

\[
\nu = 1/T_w,\quad \beta_k = \lambda
\]

\[
\delta_k = \min(k,n)\mu + \max(k-n,0)\nu
\]

\[
p_k = p_{k-1}\beta_k/\delta_k,\quad k \ge 1
\]

\[
p_0 = \left(1 + \sum_{k=1}^{\infty}\prod_{i=1}^{k}\beta_i/\delta_i\right)^{-1}
\]
Во всех точках sweep бесконечный хвост распределения усечён при epsilon = 1e-12.

Максимальные верхние оценки невключённого остатка не превышают 3.369e-14 по вероятности и 9.871e-13 по среднему числу заявок в очереди, поэтому их вклад пренебрежим в пределах учебной точности данного отчёта.

Например, при n = 1 вероятность существования очереди составляет 0.998027, а средняя длина очереди равна 8.001761; при n = 8 эти значения снижаются до 0.196074 и 0.462763.

Статусы точек остаются stationary\_truncated и прямо отражают используемую численную политику.

\noindent\textbf{Занятые операторы.} После усечения хвоста на уровне текущей численной точности по распределению суммируется среднее число занятых операторов и их загрузка.

\[
M_{\mathrm{зан}} = \sum_{k=0}^{\infty}\min(k,n)\,p_k,\quad K_{\mathrm{загр}} = M_{\mathrm{зан}}/n
\]
\begin{figure}[H]
\centering
\includegraphics[width=0.8\textwidth]{\detokenize{../figures/task1_4__busy_operators_vs_operators.png}}
\caption*{Рисунок 1.4.2. Среднее число занятых операторов при изменении числа операторов.}
\end{figure}
\begin{figure}[H]
\centering
\includegraphics[width=0.8\textwidth]{\detokenize{../figures/task1_4__operators_utilization_vs_operators.png}}
\caption*{Рисунок 1.4.3. Коэффициент загрузки операторов при изменении числа операторов.}
\end{figure}
\noindent\textbf{Очередь и уход клиентов.} Хвост распределения начиная с состояний k > n даёт вероятность существования очереди и её среднюю длину уже с учётом ухода клиентов.

\[
P_{\mathrm{оч}} = \sum_{k=n+1}^{\infty}p_k
\]

\[
L_{\mathrm{оч}} = \sum_{k=n+1}^{\infty}(k-n)\,p_k
\]
\begin{figure}[H]
\centering
\includegraphics[width=0.8\textwidth]{\detokenize{../figures/task1_4__queue_exists_vs_operators.png}}
\caption*{Рисунок 1.4.4. Вероятность наличия очереди при изменении числа операторов.}
\end{figure}
\begin{figure}[H]
\centering
\includegraphics[width=0.8\textwidth]{\detokenize{../figures/task1_4__queue_length_vs_operators.png}}
\caption*{Рисунок 1.4.5. Средняя длина очереди при изменении числа операторов.}
\end{figure}
\clearpage

\section*{Задача №2. Проектирование производственного участка.}
Имеется участок с N станками. Среднее время между наладками составляет Tc минут, среднее время наладки — Ts минут. Все потоки случайных событий считать пуассоновскими.

Построить графики от числа наладчиков: математического ожидания числа простаивающих станков; математического ожидания числа станков, ожидающих обслуживания; вероятности ожидания обслуживания; математического ожидания числа занятых наладчиков; коэффициента занятости наладчиков.

\noindent\textbf{Исходные данные:}\\
N = 35\\
Tc = 101\\
Ts = 53

\begin{figure}[H]
\centering
\includegraphics[width=0.9\textwidth]{\detokenize{assets/task2_1__scheme.png}}
\caption*{Рисунок 2.1.1. Расчетная схема.}
\end{figure}
Состояние i равно числу неисправных станков; интенсивность отказов зависит от числа работающих станков, а интенсивность восстановления ограничена числом наладчиков r. Для текущего варианта λ = 0.009901 1/мин, μ = 0.018868 1/мин.

\noindent\textbf{Распределение состояний.} Вероятности состояний строятся по числу неисправных станков i, а интенсивности отказов и восстановления меняются вместе с текущим состоянием участка.

\[
\lambda_i = (N-i)\lambda,\quad \mu_i = \min(i,r)\mu
\]

\[
p_i = p_{i-1}\lambda_{i-1}/\mu_i,\quad i=1,\ldots,N
\]

\[
p_0 = \left(1 + \sum_{i=1}^{N}\prod_{j=1}^{i}\lambda_{j-1}/\mu_j\right)^{-1}
\]
\noindent\textbf{Простаивающие и ожидающие станки.} Сначала по распределению состояний вычисляются средние числа простаивающих станков и станков, ожидающих обслуживания.

\[
M_{\mathrm{пр}} = \sum_{i=0}^{N}i\,p_i,\quad M_{\mathrm{ож}} = \sum_{i=0}^{N}\max(i-r,0)\,p_i
\]
\begin{figure}[H]
\centering
\includegraphics[width=0.8\textwidth]{\detokenize{../figures/task2_1__idle_machines_vs_repairers.png}}
\caption*{Рисунок 2.1.2. Среднее число простаивающих станков при изменении числа наладчиков.}
\end{figure}
\begin{figure}[H]
\centering
\includegraphics[width=0.8\textwidth]{\detokenize{../figures/task2_1__waiting_machines_vs_repairers.png}}
\caption*{Рисунок 2.1.3. Среднее число ожидающих станков при изменении числа наладчиков.}
\end{figure}
\noindent\textbf{Вероятность ожидания обслуживания.} Отдельно рассматривается вероятность того, что новый отказавший станок попадёт именно в ожидание, а не сразу на обслуживание.

\[
P_{\mathrm{ож}} = \frac{\sum_{i=r}^{N-1}(N-i)p_i}{\sum_{i=0}^{N-1}(N-i)p_i}
\]
\begin{figure}[H]
\centering
\includegraphics[width=0.8\textwidth]{\detokenize{../figures/task2_1__waiting_probability_vs_repairers.png}}
\caption*{Рисунок 2.1.4. Вероятность ожидания обслуживания при изменении числа наладчиков.}
\end{figure}
Вероятность ожидания в задаче 2.1 трактуется как вероятность того, что новый отказавший станок попадёт в ожидание обслуживания; это не стационарная доля состояний системы с очередью.

На концах диапазона r = 1 и r = 35 значения P\_ож составляют 1 и 0 соответственно.

Для характерной точки r = 10 получаем P\_ож = 0.918537, тогда как стационарная вероятность наличия очереди равна 0.908207; среднее число ожидающих станков при этом составляет 6.270186.

\noindent\textbf{Занятые наладчики.} После этого по тем же вероятностям состояний определяются среднее число занятых наладчиков и коэффициент их занятости.

\[
M_{\mathrm{зан}} = \sum_{i=0}^{N}\min(i,r)\,p_i,\quad K_{\mathrm{загр}} = M_{\mathrm{зан}}/r
\]
\begin{figure}[H]
\centering
\includegraphics[width=0.8\textwidth]{\detokenize{../figures/task2_1__busy_repairers_vs_repairers.png}}
\caption*{Рисунок 2.1.5. Среднее число занятых наладчиков при изменении числа наладчиков.}
\end{figure}
\begin{figure}[H]
\centering
\includegraphics[width=0.8\textwidth]{\detokenize{../figures/task2_1__repairers_utilization_vs_repairers.png}}
\caption*{Рисунок 2.1.6. Коэффициент занятости наладчиков при изменении числа наладчиков.}
\end{figure}
\clearpage

\end{document}
